% ------------------------------------------------------------------ glossary
\mwnewpart
\renewcommand\runningtitle{Glossary}
\namedpages{Glossary}
\chapter*{Glossary}
\setcounter{page}{1}
\addcontentsline{toc}{chapter}{Glossary}

\newcommand\gl[2]{\par\noindent\hangindent=1.2em\textbf{#1}\quad #2\par}

\gl{8088}{The Intel processor of the IBM Personal Computer and XT. It has
20 address lines and can address 1 MB of memory, and moves data 8 bits at
a time.}
\gl{active partition}{The partition the computer starts from.}
\gl{allocation unit}{The unit in which disk space is given to files: one or
more sectors. Also called a cluster.}
\gl{archive attribute}{A mark on a file that says it has changed since it
was last backed up.}
\gl{AUTOEXEC.BAT}{The batch file run each time the computer starts.}
\gl{batch file}{A file of commands, run by typing its name.}
\gl{BIOS}{Basic Input/Output System: the programs in the computer's ROM that
start it and work its hardware.}
\gl{boot record}{The first sector of a disk or partition, which starts the
system and describes the volume. Also called the startup record. The
first sector of a fixed disk is its master boot record, which holds the
partition table.}
\gl{buffer}{Memory in which RetrOS/LT keeps what it last read from or will
write to a disk.}
\gl{command processor}{\file{COMMAND.COM}: the program that shows the
prompt and carries out commands.}
\gl{CONFIG.SYS}{The file RetrOS/LT reads when the computer starts, to set
itself up and load device drivers.}
\gl{current directory}{The directory RetrOS/LT uses when you name none.}
\gl{current drive}{The drive RetrOS/LT uses when you name none.}
\gl{device driver}{A program loaded by a \typed{DEVICE=} line in
\file{CONFIG.SYS}, which stays in memory and works a device.}
\gl{directory}{A list of files and other directories on a disk.}
\gl{environment}{A set of named values, such as \file{PATH}, that every
program can read.}
\gl{external command}{A command that is a program file on the disk.}
\gl{file allocation table}{The table on a disk that records which
allocation units belong to which file. Under DOS 2 its entries are 12
bits long.}
\gl{fixed disk}{A disk that is not removed from the computer; a hard
disk.}
\gl{internal command}{A command that is part of the command processor.}
\gl{KB}{Kilobyte: 1,024 bytes.}
\gl{MB}{Megabyte: 1,024 KB.}
\gl{partition}{A part of a fixed disk set aside for one operating system.}
\gl{path}{The directories that lead to a file; also, the list of
directories searched for programs.}
\gl{pipe}{A bar between two commands, which makes what the first writes
into what the second reads.}
\gl{read-only}{A mark on a file that keeps it from being changed or
erased.}
\gl{redirection}{Sending what a command writes to a file or device, or
having it read from one.}
\gl{root directory}{The first directory of a disk, written as a
backslash.}
\gl{sector}{The unit in which a disk is read and written, usually 512
bytes.}
\gl{system files}{\file{RETBIO.COM} and \file{RETDOS.COM}: RetrOS/LT
itself, hidden at the start of the startup disk.}
\gl{track}{One of the circles on a disk on which data is recorded. A
5-1/4 inch diskette of the PC has 40 on each side.}
\gl{volume label}{A name of up to eleven characters given to a disk.}
\gl{wildcard}{A \typed{?} or \typed{*} in a file name, standing for other
characters.}


% ------------------------------------------------------------------ index
\mwnewpart
\renewcommand\runningtitle{Index}
\namedpages{Index}
\setcounter{page}{1}
\indexprologue{}
\printindex

% ------------------------------------------------------------------ quick reference
\mwcutpage
\pagestyle{blank}
\newgeometry{left=\dimexpr0.6in+\bleed\relax,right=\dimexpr0.6in+\bleed\relax,top=\dimexpr0.55in+\bleed\relax,bottom=\dimexpr0.6in+\bleed\relax}
\renewcommand\thepage{}
\begingroup\small\setlength\parskip{3pt}
\noindent\begin{tikzpicture}\addressband{(0,0)}{\linewidth}{5pt}\end{tikzpicture}\par
\noindent{\cond\bfseries\fontsize{20}{22}\selectfont\color{disc}RetrOS/LT Quick Reference Card}\hfill{\footnotesize Version 1.0}\par
\vspace{4pt}
\begin{minipage}[t]{0.485\linewidth}\setlength\parskip{3pt}
\textbf{\large Files}\par
\begin{tabular}{@{}P{1.7in}P{1.6in}@{}}
\typed{DIR} [\var{file}] [/P] [/W] & list files\\
\typed{COPY} \var{source} [\var{target}] & copy files\\
\typed{REN} \var{file} \var{newname} & rename\\
\typed{DEL} \var{file} & erase\\
\typed{TYPE} \var{file} & show on the screen\\
\typed{PRINT} \var{file} & print\\
\typed{COMP} \var{file1} \var{file2} & compare\\
\end{tabular}

\textbf{\large Directories}\par
\begin{tabular}{@{}P{1.7in}P{1.6in}@{}}
\typed{MD} \var{path} & make\\
\typed{CD} \var{path}\quad \typed{CD ..}\quad \typed{CD \bs} & change\\
\typed{RD} \var{path} & remove (must be empty)\\
\typed{TREE} [\var{d:}] [/F] & show the tree\\
\typed{PATH} \var{d:path}[;...] & where programs are found\\
\end{tabular}

\textbf{\large Disks}\par
\begin{tabular}{@{}P{1.7in}P{1.6in}@{}}
\typed{FORMAT B:} & format a 360 KB diskette\\
\typed{FORMAT B:/1} & format a 180 KB diskette\\
\typed{FORMAT B:/S} & make a startup diskette\\
\typed{DISKCOPY A: B:} & copy a diskette\\
\typed{CHKDSK} [\var{d:}] [/F] & check a disk\\
\typed{VOL} [\var{d:}] & volume label\\
\typed{SYS} \var{d:} & copy the system files\\
\typed{FDISK} & partition a fixed disk\\
\typed{BACKUP C:\bs\ A: /S} & back up drive C\\
\typed{RESTORE A: C:\bs*.* /S} & restore it\\
\end{tabular}
\end{minipage}\hfill
\begin{minipage}[t]{0.485\linewidth}\setlength\parskip{3pt}
\textbf{\large Wildcards and redirection}\par
\begin{tabular}{@{}P{1.2in}P{2.1in}@{}}
\typed{?}\quad \typed{*} & any one character; any characters\\
\typed{>} \var{file}\quad \typed{>>} \var{file} & write to a file; add to it\\
\typed{<} \var{file} & read from a file\\
\var{a} \typed{|} \var{b} & pipe \var{a} into \var{b}: \typed{DIR | MORE}\\
\end{tabular}

\textbf{\large Keys}\par
\begin{tabular}{@{}P{1.45in}P{1.85in}@{}}
\keys{Ctrl}{Break} & stop a command\\
\keys{Ctrl}{Num Lock} & stop the screen; any key goes on\\
\key{F3} & repeat the last command\\
\key{F1} & repeat one character of it\\
\key{Esc} & cancel the line\\
\key{Shift}\,+\,\key{PrtSc} & print the screen\\
\key{Ctrl}\,+\,\key{PrtSc} & print as you go, on or off\\
\key{Ctrl}\,+\,\key{Alt}\,+\,\key{Del} & restart the computer\\
\end{tabular}

\textbf{\large The system}\par
\begin{tabular}{@{}P{1.45in}P{1.85in}@{}}
\typed{VER} & the version\\
\typed{DATE}\quad \typed{TIME} & show or set (year in four digits)\\
\typed{SET} [\var{name}=\var{value}] & the environment\\
\typed{PROMPT \$P\$G} & show the directory\\
\typed{MODE CO80}\quad \typed{MODE MONO} & the display\\
\typed{CLS} & clear the screen\\
\end{tabular}

\textbf{\large As installed}\par
\shown{FILES=20}\quad\shown{BUFFERS=15}\quad(\shown{10} below 384 KB)\par
\shown{DEVICE=C:\bs{}DOS\bs{}ANSI.SYS}\par
\shown{DATE}, \shown{TIME}, \shown{PROMPT \$P\$G}, \shown{PATH C:\bs{}DOS}\par

\textbf{If the computer does not start}\par
Start from the DOS diskette and press \key{Esc}. For a missing system:
\typed{SYS C:}, then copy \file{COMMAND.COM} to \file{C:\bs}.
\end{minipage}
\par\vfill
\noindent{\footnotesize\color{grey}Cut along the outer edge. The Lyon Mountain Company, Republic Lyon Mountain Corporation, Lyon Mountain, New York 12952.}
\endgroup
\mwcutpage

% ------------------------------------------------------------------ reader's comment form
\begingroup\small\setlength\parskip{5pt}
\noindent{\cond\bfseries\fontsize{18}{20}\selectfont\color{disc}Reader's Comment Form}\hfill
{\footnotesize\textbf{RetrOS/LT Guide to Operations}}\par
\vspace{2pt}\noindent\begin{tikzpicture}\addressband{(0,0)}{\linewidth}{4pt}\end{tikzpicture}\par

Your comments help us improve our publications. Please tell us about any
errors, about anything that was hard to follow, and about anything you
looked for and could not find. Comments about RetrOS/LT itself are
welcome too; if a program does not run under it, please name the program
and say what happened.

\newcommand\cbox{\tikz\draw (0,0) rectangle (0.11in,0.11in);}
\begin{tabular}{@{}lll@{}}
\multicolumn{3}{@{}l}{Did you find the material:}\\
\cbox\ Easy to read and understand? & \cbox\ Yes & \cbox\ No\\
\cbox\ Organized for convenient use? & \cbox\ Yes & \cbox\ No\\
\cbox\ Complete? & \cbox\ Yes & \cbox\ No\\
\cbox\ Well illustrated? & \cbox\ Yes & \cbox\ No\\
\end{tabular}

Computer:\quad \cbox\ IBM PC\quad \cbox\ IBM PC XT\quad \cbox\ Other: \rule{1.6in}{0.4pt}\qquad BIOS date: \rule{1.0in}{0.4pt}

Memory and disks: \rule{5.2in}{0.4pt}

Comments (please give page numbers):

\foreach \n in {1,...,11}{\noindent\rule{\linewidth}{0.4pt}\par\vspace{9pt}}

\begin{tabular}{@{}p{0.9in}p{5.6in}@{}}
Name & \rule{5.4in}{0.4pt}\\[8pt]
Address & \rule{5.4in}{0.4pt}\\[8pt]
& \rule{5.4in}{0.4pt}\\
\end{tabular}

\vfill
\noindent\makebox[\linewidth]{\footnotesize\color{grey}\rule[2pt]{1.6in}{0.3pt}\ \ Fold here, with this side inside\ \ \rule[2pt]{1.6in}{0.3pt}}
\endgroup
\clearpage

\vspace*{\fill}
\begin{center}
\begin{tikzpicture}[x=1in,y=1in]
  \draw[dashed,grey] (-3.3,1.6) -- (3.3,1.6);
  \node[font=\footnotesize,text=grey] at (0,1.75) {Fold here};
  \draw (2.35,0.55) rectangle (3.05,1.35);
  \node[font=\scriptsize,align=center,text=grey] at (2.7,0.95) {Place\\stamp\\here};
  \node[align=left,anchor=north west,font=\normalsize] at (-1.2,0.2) {%
    \textbf{RetrOS/LT Comment Form}\\
    P.O. Box 290\\
    Lyon Mountain, New York 12952 USA};
  \node[anchor=north west,font=\footnotesize] at (-3.3,1.35) {From:};
  \draw[dashed,grey] (-3.3,-1.6) -- (3.3,-1.6);
  \node[font=\footnotesize,text=grey] at (0,-1.75) {Fold here, and tape closed. Please do not staple.};
\end{tikzpicture}
\end{center}
\vspace*{\fill}
\restoregeometry
